

```latex
% The following web sources are cited in the document (markdown form as requested):
% [exa.ai](https://exa.ai/library/publication/lbqpyt46khy)
% [ar5iv.labs.arxiv.org](https://ar5iv.labs.arxiv.org/html/1002.2366)
% [ar5iv.labs.arxiv.org](https://ar5iv.labs.arxiv.org/html/1111.0111)
% [jep.centre-mersenne.org](https://jep.centre-mersenne.org/articles/10.5802/jep.283/)

\documentclass[11pt]{article}

\usepackage[utf8]{inputenc}
\usepackage[T1]{fontenc}
\usepackage{amsmath,amssymb,amsthm}
\usepackage{hyperref}

\newtheorem{theorem}{Theorem}
\newtheorem{lemma}{Lemma}
\newtheorem{proposition}{Proposition}
\newtheorem{corollary}{Corollary}
\theoremstyle{definition}
\newtheorem{definition}{Definition}
\theoremstyle{remark}
\newtheorem{remark}{Remark}

\numberwithin{equation}{section}

\title{Zero topological entropy for smooth flows satisfying a distributional hypoellipticity condition}
\author{}
\date{}

\begin{document}

\maketitle

\begin{abstract}
Let \(M\) be a compact connected smooth manifold, \(X\) a smooth vector field, and
\(\phi_t\) its flow. Assume that whenever \(f\in C^\infty(M)\), \(u\in\mathcal D'(M)\),
and \(\mathcal L_X u=f\), then \(u\) is a smooth function. We prove that
\(\phi_t\) has zero topological entropy. The argument reduces the hypothesis to
global hypoellipticity of \(\mathcal L_X\), shows that positive topological entropy
would force the existence of a smooth ergodic invariant measure of positive metric
entropy, and then invokes Herman's theorem that positive metric entropy for a smooth
flow produces a non-smooth distributional solution of a smooth cohomological equation.
\end{abstract}

\section{Preliminaries and notation}

Let \(\mathcal D'(M)\) denote the space of distributions on \(M\). For the flow
\(\phi_t\) generated by \(X\), the Lie derivative on distributions is defined by
\[
  \mathcal L_X u
  = \left.\frac{d}{dt}\right|_{t=0} \phi_t^*u,
\]
where \(\phi_t^*\) is the pull-back on distributions. With this convention,
if \(u\in C^\infty(M)\), then \(\mathcal L_Xu=Xu\).

A Borel probability measure \(\mu\) is identified with a distribution of order \(0\).
It is \(\phi_t\)-invariant if and only if
\[
  \mathcal L_X\mu=0.
\]

We use the standard convention that the topological entropy of a flow is the
topological entropy of its time-one map:
\[
  h_{\mathrm{top}}(\phi_t):=h_{\mathrm{top}}(\phi_1).
\]
See, for instance, the discussion in \cite{exa} and the standard references
\cite{walters,kh}.

The hypothesis in the problem is exactly the statement that the first-order
operator \(\mathcal L_X\) is globally hypoelliptic in the following sense.

\begin{definition}
A smooth vector field \(X\) on a compact smooth manifold \(M\) is called
\emph{globally hypoelliptic} if
\[
  u\in\mathcal D'(M),\quad \mathcal L_Xu\in C^\infty(M)
  \quad\Longrightarrow\quad
  u\in C^\infty(M).
\]
\end{definition}

Thus the assumption says that \(X\) is globally hypoelliptic. This notion and its
basic consequences were studied systematically by Greenfield and Wallach
\cite{gw}; see also \cite{herman,kh}.

\section{Invariant measures under the hypoellipticity assumption}

We first record an elementary but important consequence of the hypothesis:
every invariant probability measure is in fact given by a smooth positive density.
This will allow us to apply Herman's theorem in the smooth-measure setting.

\begin{lemma}\label{lem:smooth-invariant-measures}
Assume \(X\) satisfies the hypothesis of the problem. Then every
\(\phi_t\)-invariant Borel probability measure on \(M\) is given by a
\(C^\infty\) strictly positive density. In particular, the flow is minimal and
uniquely ergodic.
\end{lemma}

\begin{proof}
Let \(\mu\) be a \(\phi_t\)-invariant Borel probability measure. As a measure,
\(\mu\) is a distribution of order \(0\), and invariance means
\[
  \mathcal L_X\mu=0.
\]
By the hypothesis, with \(f=0\), the distribution \(\mu\) must be a smooth
function. Thus, after fixing any smooth positive volume form \(\Omega\) on \(M\),
there is a nonnegative function \(h\in C^\infty(M)\) such that
\[
  d\mu=h\,d\Omega.
\]

Let \(K=\operatorname{supp}\mu\). Since \(\mu\) is invariant and each \(\phi_t\) is a
homeomorphism, \(K\) is a nonempty compact \(\phi_t\)-invariant set. Suppose
\(K\neq M\). Then \(\partial K\) is a nonempty compact \(\phi_t\)-invariant set.
By the Krylov--Bogolyubov theorem \cite{walters}, there exists a
\(\phi_t\)-invariant Borel probability measure \(\nu\) supported on \(\partial K\).
Applying the hypothesis to \(\nu\), we find that \(\nu\) is given by a smooth
density. But a nonzero smooth density cannot be supported on a closed set with
empty interior, and \(\partial K\) has empty interior. This contradiction shows
that \(K=M\). Hence every invariant probability measure has full support.

Now suppose that \(h\) vanishes at some point. Then the zero set
\[
  Z=\{x\in M:h(x)=0\}
\]
is a nonempty compact \(\phi_t\)-invariant set. Again by Krylov--Bogolyubov,
there exists an invariant probability measure supported on \(Z\). By the previous
paragraph, every invariant probability measure has full support, so \(Z=M\),
contradicting \(\mu(M)=1\). Therefore \(h>0\) everywhere.

It remains to prove minimality and unique ergodicity. If an orbit closure
\(\overline{\{\phi_t x:t\in\mathbb R\}}\) were a proper subset of \(M\), then by
Krylov--Bogolyubov it would support an invariant probability measure, whose
support would have to be all of \(M\), a contradiction. Thus every orbit is
dense, and the flow is minimal.

Finally, let \(\mu_i=h_i\Omega\), \(i=1,2\), be two invariant probability
measures. We have just shown that \(h_1,h_2\) are smooth and strictly positive.
Let \(J_t\) denote the Jacobian of \(\phi_t\) with respect to \(\Omega\). The
invariance of \(\mu_i\) gives
\[
  h_i(\phi_t x)J_t(x)=h_i(x)
  \qquad (i=1,2).
\]
Therefore the smooth function
\[
  r(x)=\frac{h_1(x)}{h_2(x)}
\]
satisfies \(r(\phi_t x)=r(x)\) for all \(t\). Since the flow is minimal and \(r\)
is continuous, \(r\) is constant. Hence \(\mu_1=\mu_2\), proving unique
ergodicity.
\end{proof}

\section{Herman's entropy--cohomology theorem}

The deep input needed is the following theorem of Herman. It says that positive
metric entropy for a smooth invariant measure forces the cohomological equation
\(\mathcal L_Xu=f\) to have a non-smooth distributional solution for some smooth
right-hand side.

\begin{theorem}[Herman]\label{thm:herman}
Let \(Y\) be a \(C^\infty\) vector field on a compact smooth manifold \(N\), and
let \(\psi_t\) be its flow. Let \(m\) be a \(C^\infty\) ergodic
\(\psi_t\)-invariant probability measure with strictly positive smooth density.
If
\[
  h_m(\psi_1)>0,
\]
then there exist
\[
  f\in C^\infty(N), \qquad u\in\mathcal D'(N)\setminus C^\infty(N)
\]
such that
\[
  \mathcal L_Yu=f.
\]
Equivalently, a globally hypoelliptic smooth vector field cannot have positive
metric entropy with respect to a smooth ergodic invariant measure.
\end{theorem}

This theorem is due to Herman \cite{herman}; see also the exposition in
Katok--Hasselblatt \cite{kh}. The proof combines the variational principle,
Ruelle's inequality and the Pesin entropy formula
\cite{katok,pesinweb}, and Yomdin's volume-growth theorem \cite{yomdin}.
Roughly speaking, positive metric entropy for a smooth measure forces nonzero
Lyapunov exponents. Pesin theory then produces unstable manifolds with
exponential volume growth, and Herman constructs from these objects a
distributional solution of a smooth cohomological equation whose wave front set
is nonempty; in particular, the solution is not \(C^\infty\).

For our purposes, we need only the contrapositive form:

\begin{corollary}\label{cor:herman-gh}
If \(Y\) is globally hypoelliptic and preserves a \(C^\infty\) ergodic
probability measure with positive density, then the metric entropy of \(Y\) with
respect to that measure is zero.
\end{corollary}

\section{Proof of the main statement}

We can now prove the assertion in the problem.

\begin{theorem}
Let \(M\) be a compact connected smooth manifold, \(X\in\mathfrak X(M)\) a smooth
vector field, and \(\phi_t\) its flow. Assume that for every
\(f\in C^\infty(M)\) and every \(u\in\mathcal D'(M)\),
\[
  \mathcal L_Xu=f
  \quad\Longrightarrow\quad
  u\in C^\infty(M).
\]
Then
\[
  h_{\mathrm{top}}(\phi_t)=0.
\]
\end{theorem}

\begin{proof}
By definition, the hypothesis says precisely that \(X\) is globally
hypoelliptic.

Suppose, for contradiction, that
\[
  h_{\mathrm{top}}(\phi_t)=h_{\mathrm{top}}(\phi_1)>0.
\]
By the variational principle for topological entropy of flows
\cite{walters,exa}, there exists a \(\phi_t\)-invariant Borel probability measure
\(\mu\) such that
\[
  h_\mu(\phi_1)>0.
\]
By Lemma \ref{lem:smooth-invariant-measures}, \(\mu\) is given by a smooth
strictly positive density.

Passing to an ergodic component of \(\mu\), and using the affinity of metric
entropy in the ergodic decomposition \cite{walters}, we obtain a
\(\phi_t\)-invariant ergodic Borel probability measure \(m\) such that
\[
  h_m(\phi_1)>0.
\]
Again by Lemma \ref{lem:smooth-invariant-measures}, \(m\) has a smooth strictly
positive density.

Now apply Herman's theorem, Theorem \ref{thm:herman}, to \(Y=X\) and \(m\). It
provides \(f\in C^\infty(M)\) and \(u\in\mathcal D'(M)\setminus C^\infty(M)\)
such that
\[
  \mathcal L_Xu=f.
\]
This contradicts the hypothesis, which asserts that every distributional
solution of such an equation with smooth right-hand side must itself be a smooth
function.

Therefore the assumption \(h_{\mathrm{top}}(\phi_t)>0\) is impossible. Hence
\[
  h_{\mathrm{top}}(\phi_t)=0.
\]
\end{proof}

\begin{remark}
The conclusion is consistent with the general phenomenon that zero topological
entropy is robust under many orbit equivalences and time changes, while finite
nonzero entropy is not; see, for example, \cite{exa,equiv}. Zero entropy smooth
flows can nevertheless have rich statistical behavior; for instance, smooth zero
entropy flows satisfying a classical central limit theorem are constructed in
\cite{jep}.
\end{remark}

\begin{thebibliography}{99}

\bibitem{walters}
P. Walters,
\emph{An Introduction to Ergodic Theory},
Graduate Texts in Mathematics, Vol. 79, Springer-Verlag, 1982.

\bibitem{gw}
S. Greenfield and N. Wallach,
\emph{Global hypoellipticity and Liouville numbers},
Trans. Amer. Math. Soc. 197 (1973), 211--228.

\bibitem{herman}
M. R. Herman,
\emph{Sur les courbes invariantes par les diff\'eomorphismes de l'anneau},
Ast\'erisque 62--63, Soci\'et\'e Math\'ematique de France, 1979.
See in particular the theorem on globally hypoelliptic vector fields and
entropy.

\bibitem{kh}
A. Katok and B. Hasselblatt,
\emph{Introduction to the Modern Theory of Dynamical Systems},
Encyclopedia of Mathematics and its Applications, Vol. 54,
Cambridge University Press, 1995.

\bibitem{katok}
A. Katok,
\emph{Lyapunov exponents, entropy and periodic orbits for diffeomorphisms},
Publ. Math. Inst. Hautes \'Etudes Sci. 49 (1979), 137--173.

\bibitem{yomdin}
Y. Yomdin,
\emph{Volume growth and entropy},
Israel J. Math. 57 (1987), 285--300.

\bibitem{pesinweb}
On the entropy of conservative flows,
arXiv:1002.2366,
\url{https://ar5iv.labs.arxiv.org/html/1002.2366}.

\bibitem{exa}
Topological entropies of equivalent smooth flows,
\url{https://exa.ai/library/publication/lbqpyt46khy}.

\bibitem{equiv}
Entropy and periodic orbits for equivalent smooth flows,
arXiv:1111.0111,
\url{https://ar5iv.labs.arxiv.org/html/1111.0111}.

\bibitem{jep}
D. Dolgopyat and collaborators,
\emph{Smooth zero entropy flows satisfying the classical central limit theorem},
Journal de l'\'Ecole polytechnique --- Math\'ematiques,
\url{https://jep.centre-mersenne.org/articles/10.5802/jep.283/}.

\end{thebibliography}

\end{document}
```